%% Cover Letter — Microsoft Research Cambridge, Senior Machine Learning Research Engineer
%% Compile: lualatex cover_msr-cambridge_ml-research-engineer.tex
\documentclass[11pt,a4paper]{moderncv}
\moderncvstyle{banking}
\moderncvcolor{blue}
\usepackage[scale=0.80]{geometry}
\usepackage{hyperref}

\name{Oluwatosin}{Ayeni}
\address{Glasgow, Scotland, UK}{}{}
\phone[mobile]{07727305230}
\email{ayenioluwatosinolawale@gmail.com}
\extrainfo{%
  \href{https://github.com/AyeniOluwatosinOlawale}{github.com/AyeniOluwatosinOlawale} \quad
  \href{https://www.linkedin.com/in/oluwatosinolawaleayeni/}{linkedin.com/in/oluwatosinolawaleayeni}
}

\begin{document}
\makecvtitle

\today

\vspace{1em}
Hiring Team\\
Research Engineering, Microsoft Research Cambridge\\
Cambridge, UK

\vspace{1em}
\textbf{Re: Senior Machine Learning Research Engineer — Research Engineering, MSR Cambridge}

\vspace{1em}
Dear Hiring Team,

I am writing to apply for the Senior Machine Learning Research Engineer position at Microsoft Research Cambridge. MSR's work at the intersection of foundational research and large-scale systems is exactly the environment I am seeking — a place where rigorous benchmarking, novel training approaches, and production deployment are all part of the same conversation.

My experience spans the full ML pipeline at GPU scale: from tokeniser and architecture design to training loop, evaluation, and production inference deployment. At Nexora Academic I conducted a head-to-head benchmark of vLLM v0.30 and SGLang v0.5.20 on 2$\times$ NVIDIA H100 NVL 94\,GB under TP=2 — producing 576+ measurement records across 9 context lengths and 8 concurrency levels. The headline result was a 1,443$\times$ TTFT speedup (39\,ms vs 55,565\,ms) via SGLang's RadixAttention prefix caching at 131K-token context, with 100\% SLO compliance at concurrency 1--128. I built the full benchmarking harness that produced these results: automated metric collection (TTFT, TPOT, throughput, SLO), live dashboard, and comparative reporting across closed and open serving engines.

On the training side, I have fine-tuned Mistral 7B on a radiology corpus using LoRA/PEFT — managing the complete pipeline from dataset curation and tokenisation through training, evaluation, and production API deployment, achieving +22\% diagnostic accuracy and $-$35\% hallucination rate. I subsequently fine-tuned Phi-4 on the Cosmopedia dataset with LoRA and trained a Qwen3-based language model from scratch — tokeniser, architecture configuration, training loop, and generation evaluation. These are not isolated experiments; each one ended with a deployable artefact and a set of evaluation metrics I could defend.

My GPU systems work sits alongside the modelling: CUDA KV cache profiling with NVIDIA Nsight that identified a memory bandwidth bottleneck and yielded a 3.8$\times$ throughput improvement; continuous batching pipeline redesign that reduced inference latency 18$\times$; and a 28$\times$ LLM cost reduction achieved by deploying Qwen2.5-7B on RTX 4090 hardware via vLLM/SGLang. I understand what happens between the Python API call and the GPU kernel, and I can profile and fix it at both levels.

What draws me to MSR Cambridge specifically is the combination of research depth and engineering rigour. The work you publish on model training, evaluation methodology, and inference systems is work I engage with directly — not as background reading but as reference material for decisions I am making week to week. I am open to relocating to Cambridge.

I would welcome the opportunity to discuss how my experience aligns with the team's current research engineering priorities.

Yours sincerely,

\textbf{Oluwatosin Ayeni}\\
\href{mailto:ayenioluwatosinolawale@gmail.com}{ayenioluwatosinolawale@gmail.com}\\
07727305230

\end{document}
